\documentclass[11pt, letterpaper]{article}

%% Pacotes de Idioma e Codificação
\usepackage[brazilian]{babel}
\usepackage[utf8]{inputenc}
\usepackage[T1]{fontenc}

%% Pacotes de Layout e Fontes
\usepackage[top=28mm, bottom=28mm, left=15mm, right=15mm]{geometry}
\usepackage{fontspec}
\setmainfont{Arial}
\setlength{\parindent}{2em}
\renewcommand{\baselinestretch}{1.5}

%% Pacotes Matemáticos e Técnicos
\usepackage{amsmath, amsthm, amssymb, amsfonts}
\usepackage{mathtools, bm, dsfont}
\usepackage{algorithm} 
\usepackage{algpseudocode} 

%% Pacotes de Elementos Visuais
\usepackage{graphicx}
\usepackage{booktabs, array, multirow}
\usepackage{caption}
\usepackage[section]{placeins} % Impede que figuras flutuem para outras seções
\usepackage[many]{tcolorbox}
\usepackage{xcolor}

%% Referências e Hiperlinks
\usepackage[colorlinks=true, allcolors=blue]{hyperref}
\usepackage{url}
\usepackage{natbib}

\hypersetup{
    citecolor=black,
    linkcolor=black,
    urlcolor=black
}

\begin{document}

%%%%%%%%%%%%%%%%%%%%%%%%%%%%%%%%%%%%%% 1.1 - PÁGINA DE CAPA %%%%%%%%%%%%%%%%%%%%%%%%%%%%%%%%%%%%%%%%%%%%%%%%%

\thispagestyle{empty}
\begin{center}
    
    \vspace*{-1.5cm}
    \hspace*{-0.1cm} {\rule[-1ex]{16cm}{0.01cm}}
    \begin{center}
    \hspace{0.2cm}
    \begin{minipage}[s]{2.8cm}
    % Certifique-se de que os arquivos de imagem existam na pasta Figuras/
    \includegraphics[width=2.8cm]{Figuras/logounicamp.eps} 
    \end{minipage}
    \begin{minipage}[s]{11cm}
    {
        \begin{center} 
        {\LARGE Universidade Estadual de Campinas}\\ \vspace{0.2 cm}
        {Instituto de Matemática, Estatística e Computação Científica }\\ \vspace{0.1 cm}
        {Departamento de Matemática Aplicada}
        \end{center}
    }
    \end{minipage}
    \begin{minipage}[s]{3cm}
        \includegraphics[width=2.3cm]{Figuras/logoimecc.eps}
    \end{minipage}
    \end{center}

\vspace{2 cm}

% Nome da Disciplina (Exigência 1.1)
\large{\textbf{Nome da Disciplina: [Inserir Nome da Disciplina]}} \\

\vspace{1 cm}

%% Título do Trabalho
\LARGE{
\textbf{Solução Problemas de Programação Matemática \\ 
de Grande Porte via Métodos de Pontos Interiores}
\\}

\vspace{2.5cm}

% Informações do Aluno (Exigência 1.1)
\large
{Aluno: Aluno A \hspace{0.5cm} RA: 5879671 }\\
\vspace{0.5cm}
{Data de Entrega: \today} % Altere para a data específica se necessário

\vspace{1.5cm}

\begin{abstract}
\normalsize
\textbf{Resumo das Técnicas e Resultados:} Utilizando os métodos de pontos interiores para otimização linear como solução para resolver grandes problemas que envolvem matrizes esparsas, analisamos a redução do tempo computacional por iteração entre os projetos PCx0 e PCx13, reduzindo o número de iterações para alcançar a convergência de métodos e o desenvolvimento de métodos para problemas específicos com estruturas particulares.
\end{abstract}

\end{center}

\newpage
\pagenumbering{arabic}

%%%%%%%%%%%%%%%%%%%%%%%%%% 1.2 - PÁGINAS SUBSEQUENTES %%%%%%%%%%%%%%%%%%%%%%%%%%

\section{Fundamentação Teórica e Literatura}
Apresente aqui a revisão bibliográfica e a literatura estudada para a implementação da tarefa. Cite os autores e métodos relevantes utilizando o formato adequado.

\section{Implementação e Metodologia}
Descreva aqui os detalhes técnicos da execução do projeto.

\subsection{Algoritmos Implementados}
Utilize o ambiente de algoritmos para descrever sua lógica:

\begin{algorithm}
\caption{Método de Pontos Interiores (Exemplo)}
\begin{algorithmic}[1]
\State Inicializar $x^0, s^0 > 0$
\While{Critério de parada não satisfeito}
    \State Calcular direções de busca $\Delta x, \Delta s$
    \State Atualizar variáveis com passo $\alpha$
\EndWhile
\end{algorithmic}
\end{algorithm}

\subsection{Dificuldades Encontradas e Soluções Propostas}
Relate os obstáculos técnicos ou teóricos e como foram superados durante o desenvolvimento.

\section{Resultados e Discussão}
Apresente aqui os resultados quantitativos e qualitativos.

\begin{table}[h]
\centering
\caption{Comparação de performance entre PCx0 e PCx13.}
\begin{tabular}{@{}lll@{}}
\toprule
Métrica & PCx0 & PCx13 \\ \midrule
Tempo Médio (s) & 150.5 & 85.2 \\
Iterações & 45 & 28 \\ \bottomrule
\end{tabular}
\end{table}

\begin{figure}[h]
\centering
%\includegraphics[width=0.6\textwidth]{Figuras/grafico_convergencia.png}
\caption{Discussão sobre a convergência dos métodos implementados.}
\end{figure}

%%%%%%%%%%%%%%%%%%%%%%%%%%%%%% 1.3 - PÁGINA FINAL %%%%%%%%%%%%%%%%%%%%%%%%%%%%%%

\newpage
\section{Conclusão}
Apresente as considerações finais sobre o que foi realizado na tarefa, destacando se os objetivos iniciais foram atingidos e a relevância dos resultados obtidos para problemas de grande porte.

\section{Sugestões para Aprimoramento}
Enumere possíveis melhorias para trabalhos futuros, como otimização de novos parâmetros, uso de diferentes estruturas de dados para matrizes esparsas ou integração com outras técnicas de deep learning.

\newpage
\bibliographystyle{plain}
\bibliography{referencias} % Certifique-se de ter um arquivo referencias.bib

\end{document}